\section{Related Work}
\label{sec:related_works}


\subsection{Spatial Intelligence Benchmark}
Multiple spatial intelligence benchmarks~\cite{Yang_2025_CVPR,yu2025far,liu2025mirage,gong2025space,deng2025internspatial,yang2025mmsi,ma20243dsrbench,cheng2024spatialrgpt,zhao2025mmvu,tang2025lego} have recently emerged to evaluate MLLM's capabilities in spatial reasoning tasks.
VSI-Bench utilizes~\cite{Yang_2025_CVPR} public 3D scene datasets including ScanNet~\cite{dai2017scannet}, ScanNet++~\cite{yeshwanth2023scannet++} and ARKitScenes~\cite{baruch2021arkitscenes} to construct 5,000 QA pairs that span eight 3D spatial cognition tasks, categorized into three types of configurational, measurement estimation and spatiotemporal reasoning.
VLM4D~\cite{zhou2025vlm4d} uses both real-world clips and synthetic videos to produce over 2,000 high-quality QA pairs to explore models' spatiotemporal abilities in translational and rotational motion, perspective shifts, motion continuity, and related dynamics.
STI-Bench~\cite{li2025sti} is built from over 300 real-world videos with more than 2,000 QA pairs to test both static and dynamic spatial tasks.
More recently, VSI-SUPER~\cite{yang2025cambrian} complements VIS-Bench by adding long videos in the benchmark to construct a large instruction-tuning dataset named VSI-590K.
While these benchmarks emphasize various spatial capabilities in different scenarios, there is a lack of a comprehensive benchmark that unifies multi-task spatial understanding and reasoning in the context of a unified 4D representation.
Unlike existing benchmarks, Spatial4D-Bench is a large-scale, multi-task evaluation benchmark comprising \textasciitilde40,000 question-answer pairs covering 18 well-defined tasks which are systematically organized into 6 categories.
Spatial4D-Bench also presents seven tasks that are important for spatial intelligence but have not been sufficiently investigated in current benchmarks.
The statistical differences between Spatial4D-Bench and existing benchmarks are summarized in \cref{tab:recent_mllm_eval_benchmarks}.



\subsection{Multimodal Large Language Models}
Multimodal Large Language Models extend language models~\cite{gemini2023,llama2_2023} by integrating extra input modalities such as video, image, and audio, to enable advanced spatial capabilities including visual reasoning and scene understanding with temporal information~\cite{Alayrac2022Flamingo,chen2024internvl}.
Qwen-VL~\cite{bai2025qwen2,qwen3} focuses on architectural design to improve high-resolution visual recognition and OCR/text understanding inside images.
InternVL~\cite{wang2025internvl3,chen2024internvl} emphasizes scalable vision encoders and multimodal preference optimization to enable scalable training pipelines.
LLaVA~\cite{liu2023improvedllava} introduces vision instruction tuning to teach the model to follow complex multimodal instructions and makes it possible to train high-quality vision language models with relatively little manually curated multimodal data.
VideoLLama~\cite{zhang2025videollama} balances input flexibility and model efficiency to support long-video QA, temporal reasoning, temporal grounding.
These improvements jointly impel MLLMs to be more capable of achieving a wide variety of challenging spatial tasks.



\subsection{Spatial Reasoning with MLLMs}
In addition to benchmarking MLLMs, some researchers have also proposed various techniques to enhance the spatial reasoning capabilities of MLLMs.
SpatialRGPT~\cite{cheng2024spatialrgpt} develops automated data labeling pipelines to generate large quantities of 3D visual QA pairs for training, with the integration of depth modules and the alignment of depth information with 2D visual embeddings.
Tang \etal~\cite{tang2025sparkle} investigates whether enhancing basic spatial abilities such as direction comprehension, distance estimation, and localization can improve the overall spatial reasoning performance.
Chen \etal~\cite{chen2025spatial} reveals that successful spatial reasoning largely depends on the model's ability to attend to task-relevant objects.
Some other work~\cite{fan2025vlm,wu2025spatial,zhu2024llava} further lifts 2D encodings to 3D to recover implicit 3D structural information.
While various efforts have been made to improve the spatial reasoning capabilities of MLLMs, thorough evaluation of the model performance on a comprehensive benchmark will provide insights and understanding of how the current limitations can be further addressed.

